% =============================================================================
%                    WEEK 8: LLM TRAINING STRATEGIES
% =============================================================================
\section{Week 8: LLM Training Strategies}

% -----------------------------------------------------------------------------
%                              8.1 TOPIC SUMMARY
% -----------------------------------------------------------------------------
\subsection{Topic Summary}

\begin{keybox}[Learning Objectives]
By the end of this week, you will be able to:
\begin{enumerate}
  \item Describe the principles of \textbf{supervised fine-tuning (SFT)} and \textbf{reinforcement learning from human feedback (RLHF)}.
  \item Explain the three-stage LLM training pipeline: \textbf{pre-training} $\to$ \textbf{instruction fine-tuning} $\to$ \textbf{RLHF}.
  \item Describe the difference between \textbf{PPO}, \textbf{DPO}, and \textbf{GRPO}.
  \item Explain \textbf{LoRA} and how it reduces memory requirements for fine-tuning.
  \item Calculate PPO objectives given probability ratios and advantages.
\end{enumerate}
\end{keybox}

\separator

\subsubsection*{The Big Picture}

This week explains how LLMs are trained through pre-training, instruction fine-tuning, and human-feedback methods (RLHF, PPO, DPO, GRPO) to turn raw next-token predictors into aligned, instruction-following models. We also address challenges such as memorisation, catastrophic forgetting, annotation cost, and training stability.

\separator

\subsubsection*{What Examiners Like to Ask}

\begin{itemize}
  \item Describe the three stages of LLM training and their purposes.
  \item Explain the \textbf{loss function} and \textbf{data requirements} for pre-training vs.\ SFT vs.\ RLHF.
  \item Define \textbf{catastrophic forgetting} and other SFT challenges.
  \item Explain \textbf{RLHF} and the role of \textbf{reward models}.
  \item Calculate \textbf{PPO objectives}: ratio, clipped ratio, unclipped term, final objective.
  \item Explain \textbf{LoRA}: how it reduces parameters via matrix decomposition.
  \item Compare \textbf{DPO} and \textbf{GRPO} to PPO-based RLHF.
\end{itemize}

\bigseparator

% -----------------------------------------------------------------------------
%                 8.2 OVERVIEW + CORE CONCEPTS + EXTENDED THEORY
% -----------------------------------------------------------------------------
\subsection{Overview + Core Concepts + Extended Theory}

\subsubsection*{Roadmap}
\begin{enumerate}
  \item Challenges in LLM training
  \item The three-step training pipeline
  \item Pre-training: next-token prediction
  \item Memorisation vs.\ learning
  \item Instruction fine-tuning (SFT)
  \item SFT challenges and data considerations
  \item RLHF and reward models
  \item PPO: clipping and stability
  \item LoRA: efficient fine-tuning
  \item DPO and GRPO: simpler alternatives
\end{enumerate}

\bigseparator

% --- 8.2.1 Challenges in LLM Training ---
\subsubsection{Challenges in LLM Training}

\begin{defbox}[LLM Training as Supervised Learning?]
Traditional supervised learning has three components:
\begin{itemize}
  \item \textbf{Loss function:} The learning target
  \item \textbf{Optimiser:} How to find the target
  \item \textbf{Labelled data:} Where to train the model
\end{itemize}

\textbf{LLM challenges:}
\begin{itemize}
  \item \textbf{Loss function:} How to define correctness for language generation?
  \item \textbf{Optimiser:} How to optimise when loss is unclear?
  \item \textbf{Labelled data:} Impossible to manually label all human knowledge
\end{itemize}
\end{defbox}

\begin{tipbox}[Solution: Step-by-Step Training]
LLM training is a systematic task with no simple solution. The key insight: break it into steps, each solving part of the problem.
\end{tipbox}

\bigseparator

% --- 8.2.2 Three-Step Training Pipeline ---
\subsubsection{The Three-Step LLM Training Pipeline}

\begin{thmbox}[LLM Training Pipeline]
\begin{enumerate}
  \item \textbf{Step 1: Pre-training} --- ``Speak like a human''
  \begin{itemize}
    \item Train to predict next token
    \item Uses plain text (no labels needed)
    \item Gives fluency and world knowledge
  \end{itemize}
  
  \item \textbf{Step 2: Instruction Fine-Tuning (SFT)} --- ``Understand instructions''
  \begin{itemize}
    \item Train on instruction-response pairs
    \item Teaches task-following behavior
  \end{itemize}
  
  \item \textbf{Step 3: RLHF} --- ``Learn from interaction''
  \begin{itemize}
    \item Align with human preferences
    \item Uses comparative feedback (which response is better?)
  \end{itemize}
\end{enumerate}
\end{thmbox}

\bigseparator

% --- 8.2.3 Pre-training ---
\subsubsection{Pre-training: Next-Token Prediction}

\begin{defbox}[Pre-training]
\textbf{Goal:} Train the model to predict the next token given context.

\textbf{How it simplifies the three challenges:}
\begin{itemize}
  \item \textbf{Loss function:} Cross-entropy loss on next-token prediction
  \item \textbf{Optimiser:} Standard gradient descent
  \item \textbf{Data:} Plain text --- no annotation required!
\end{itemize}
\end{defbox}

\begin{exbox}[Pre-training Example]
\textbf{Training sentence:} ``May the force be with you''

\textbf{Training steps:}
\begin{center}
\begin{tabular}{@{} l l l @{}}
\toprule
\textbf{Context} & \textbf{Predicted} & \textbf{Target} \\
\midrule
``May the force'' & LLM $\to$ ? & ``be'' \\
``May the force be'' & LLM $\to$ ? & ``with'' \\
``May the force be with'' & LLM $\to$ ? & ``you'' \\
\bottomrule
\end{tabular}
\end{center}

\textbf{Loss:} $\mathcal{L} = -\log P(\text{target} \mid \text{context})$

If the model predicts the correct word with high probability, loss is low $\to$ small parameter updates.
\end{exbox}

\begin{tipbox}[Teacher Forcing]
During training, even if the model predicts incorrectly, we use the \textbf{correct word} as input for the next step (not the model's wrong prediction).
\end{tipbox}

\bigseparator

% --- 8.2.4 Memorisation vs Learning ---
\subsubsection{Memorisation vs.\ Learning}

\begin{tipbox}[The Risk of Complex Models]
\textbf{Trade-off:} Model complexity vs.\ generalisability.

LLMs are extremely complex (billions of parameters). This creates risks:
\begin{itemize}
  \item \textbf{Overfitting:} Learning noise in training data
  \item \textbf{Memorisation:} Storing patterns rather than understanding them
\end{itemize}

\textbf{Question:} Does the model truly understand, or just memorise?
\end{tipbox}

\begin{exbox}[Spurious Correlations Example]
\textbf{Task:} Classify Pok\'{e}mon vs.\ Digimon images.

\textbf{Model accuracy:} 98\% on training and test data.

\textbf{Reality:} The model learned that:
\begin{itemize}
  \item Pok\'{e}mon images are PNG (transparent $\to$ black background)
  \item Digimon images are JPEG (white background)
\end{itemize}

The model discriminated based on \textbf{background colour}, not the actual characters!
\end{exbox}

\bigseparator

% --- 8.2.5 Instruction Fine-Tuning ---
\subsubsection{Instruction Fine-Tuning (SFT)}

\begin{defbox}[Instruction Fine-Tuning]
\textbf{Goal:} Teach the model to follow human instructions.

\textbf{Process:}
\begin{enumerate}
  \item Collect instruction-response pairs
  \item Initialise with pre-trained parameters
  \item Fine-tune on the instruction data
\end{enumerate}

\textbf{Key difference from pre-training:}
\begin{itemize}
  \item Pre-training: Predict every next word
  \item SFT: Only train on the \textbf{response} part (instruction is just context)
\end{itemize}
\end{defbox}

\begin{exbox}[SFT Example]
\textbf{Instruction:} ``What's the iconic expression from the Star Wars movie?''

\textbf{Training (fine-tune on response only):}
\begin{center}
\begin{tabular}{@{} l l @{}}
\toprule
\textbf{Context} & \textbf{Target} \\
\midrule
{[Instruction]} & ``May'' \\
{[Instruction]} May & ``the'' \\
{[Instruction]} May the & ``force'' \\
\ldots & \ldots \\
\bottomrule
\end{tabular}
\end{center}
\end{exbox}

\begin{tipbox}[Cross-lingual Transfer]
Fine-tuning on one language can transfer to others! Example: Pre-train on 104 languages, fine-tune only on Chinese QA $\to$ model can answer English QA questions too.
\end{tipbox}

\bigseparator

% --- 8.2.6 SFT Challenges ---
\subsubsection{SFT Challenges}

\begin{defbox}[Common SFT Problems]
\begin{itemize}
  \item \textbf{Catastrophic Forgetting:} Learning new tasks overwrites previous knowledge
  \item \textbf{Instruction Ambiguity:} Different tasks expect different formats (bullets vs.\ long-form)
  \item \textbf{Data Imbalance:} Model overfits to tasks with more data
  \item \textbf{Domain Shift:} Mixing domains (medical vs.\ casual) confuses the model
\end{itemize}
\end{defbox}

\begin{tipbox}[Two Paths to Multi-Task AI]
\textbf{Path 1:} Train separate models for each task (more accurate, less general)

\textbf{Path 2:} Continuously fine-tune one model on all tasks (more general, risk of forgetting)

Most commercial LLMs use Path 2 with careful data mixing and curriculum.
\end{tipbox}

\begin{tipbox}[Data Quality Matters]
\textbf{LIMA paper:} Just 1,000 high-quality examples can match GPT-4 in 43\% of cases.

``Less Is More for Alignment'' --- quality beats quantity.
\end{tipbox}

\bigseparator

% --- 8.2.7 Self-Instruct ---
\subsubsection{Self-Instruct: Generating Training Data}

\begin{defbox}[Self-Instruct]
\textbf{Problem:} Manual annotation is expensive.

\textbf{Solution:} Use LLMs to generate instruction-response pairs automatically.

\textbf{Benefits:}
\begin{itemize}
  \item Reduces annotation cost
  \item Expands datasets beyond human-written examples
\end{itemize}

\textbf{Risks:}
\begin{itemize}
  \item Generated labels may be incorrect
  \item May violate API terms of service (e.g., OpenAI prohibits using outputs to train competing models)
\end{itemize}
\end{defbox}

\bigseparator

% --- 8.2.8 RLHF ---
\subsubsection{Reinforcement Learning from Human Feedback (RLHF)}

\begin{defbox}[RLHF]
\textbf{Problem with SFT:} Requires word-level annotation (expensive); ignores overall response quality.

\textbf{RLHF solution:} Annotators only need to \textbf{compare} responses (which is better?), not write correct answers.

\textbf{Process:}
\begin{enumerate}
  \item Generate multiple responses to a prompt
  \item Human ranks responses (or picks the better one)
  \item Increase probability of preferred responses; decrease probability of dispreferred ones
\end{enumerate}
\end{defbox}

\begin{exbox}[RLHF Example]
\textbf{Prompt:} ``What's the iconic Star Wars expression?''

\textbf{Response A:} ``May the force be with you'' $\leftarrow$ Human prefers

\textbf{Response B:} ``If necessary, please learn to let it go''

\textbf{Training:} Increase $P(\text{Response A})$, decrease $P(\text{Response B})$.
\end{exbox}

\bigseparator

% --- 8.2.9 Reward Models ---
\subsubsection{Reward Models}

\begin{defbox}[Reward Model]
\textbf{Problem:} Human annotation doesn't scale; humans struggle to distinguish good responses from well-trained models.

\textbf{Solution:} Train a \textbf{reward model} to simulate human preferences.

\textbf{Training:} Use human-ranked examples to train a model that predicts preference scores.

\textbf{Risk:} The LLM may learn to \textbf{exploit} the reward model (reward hacking), producing responses that score high but aren't actually good.
\end{defbox}

\bigseparator

% --- 8.2.10 PPO ---
\subsubsection{PPO: Proximal Policy Optimization}

\begin{defbox}[PPO Principles]
\textbf{Goal:} Update model to increase probability of good responses.

\textbf{Key principles:}
\begin{enumerate}
  \item If reward is high $\to$ increase generative probability
  \item If reward is low $\to$ decrease generative probability
  \item Don't move too far from the previous model (stability)
\end{enumerate}

\textbf{Challenge:} Requires keeping two models in memory (current + previous version).
\end{defbox}

\begin{thmbox}[PPO Clipping Mechanism]
\textbf{Ratio:} $r = \frac{P_{\text{current}}}{P_{\text{previous}}}$

\textbf{Clipped ratio:} Constrain $r$ to range $[1-\varepsilon, 1+\varepsilon]$ (typically $\varepsilon = 0.2$)

\textbf{PPO objective:}
\[
L^{\text{CLIP}} = \min\left( r \cdot A, \text{clip}(r, 1-\varepsilon, 1+\varepsilon) \cdot A \right)
\]

where $A$ is the \textbf{advantage} (reward signal).

\textbf{Why clip?} Prevents extremely large updates that could destabilise training.
\end{thmbox}

\begin{exbox}[PPO Calculation Example 1]
\textbf{Given:} $P_{\text{current}} = 0.3$, $P_{\text{previous}} = 0.2$, Advantage $= 2.0$, $\varepsilon = 0.2$

\textbf{Step 1:} Ratio $= 0.3 / 0.2 = 1.5$

\textbf{Step 2:} Unclipped term $= 1.5 \times 2.0 = 3.0$

\textbf{Step 3:} Clipped ratio $= \min(1.5, 1.2) = 1.2$ (since $1+\varepsilon = 1.2$)

\textbf{Step 4:} Clipped term $= 1.2 \times 2.0 = 2.4$

\textbf{Step 5:} Final objective $= \min(3.0, 2.4) = \boxed{2.4}$
\end{exbox}

\begin{exbox}[PPO Calculation Example 2]
\textbf{Given:} $P_{\text{current}} = 0.25$, $P_{\text{previous}} = 0.20$, Advantage $= 1.5$, $\varepsilon = 0.1$

\textbf{Step 1:} Ratio $= 0.25 / 0.20 = 1.25$

\textbf{Step 2:} Unclipped term $= 1.25 \times 1.5 = 1.875$

\textbf{Step 3:} Clipped ratio $= \min(1.25, 1.1) = 1.1$

\textbf{Step 4:} Clipped term $= 1.1 \times 1.5 = 1.65$

\textbf{Step 5:} Final objective $= \min(1.875, 1.65) = \boxed{1.65}$
\end{exbox}

\begin{tipbox}[Interpreting the Ratio]
\begin{itemize}
  \item $r > 1$: Current model already favours this response $\to$ less update needed
  \item $r < 1$: Current model doesn't favour this response $\to$ more update needed
\end{itemize}

Clipping prevents the model from over-focusing on single examples.
\end{tipbox}

\bigseparator

% --- 8.2.11 LoRA ---
\subsubsection{LoRA: Low-Rank Adaptation}

\begin{defbox}[LoRA]
\textbf{Problem:} PPO requires two models in memory $\to$ very expensive for large LLMs.

\textbf{Solution:} Instead of updating the full weight matrix $W$ (size $m \times n$), decompose updates into two smaller matrices:
\[
W' = W + L \times R
\]
where $L$ is $m \times r$ and $R$ is $r \times n$, with $r \ll m, n$.

\textbf{Parameter savings:}
\begin{itemize}
  \item Full matrix: $m \times n$ parameters
  \item LoRA: $m \times r + r \times n$ parameters
\end{itemize}
\end{defbox}

\begin{exbox}[LoRA Parameter Calculation]
\textbf{Full matrix:} $6 \times 5 = 30$ parameters

\textbf{LoRA with $r = 2$:}
\begin{itemize}
  \item $L$: $6 \times 2 = 12$ parameters
  \item $R$: $2 \times 5 = 10$ parameters
  \item Total: $12 + 10 = 22$ parameters
\end{itemize}

\textbf{Savings:} $30 - 22 = 8$ parameters (27\% reduction)

For large matrices, savings are dramatic: $1000 \times 1000 = 1M$ vs.\ $1000 \times 16 + 16 \times 1000 = 32K$.
\end{exbox}

\begin{tipbox}[LoRA Trade-offs]
\textbf{Pros:}
\begin{itemize}
  \item Reduces trainable parameters significantly
  \item Enables fine-tuning large LLMs with limited GPU memory
\end{itemize}

\textbf{Cons:}
\begin{itemize}
  \item Does NOT guarantee same accuracy as full fine-tuning
  \item Smaller $r$ $\to$ more compression but lower accuracy
  \item No free lunch: accuracy vs.\ efficiency trade-off
\end{itemize}
\end{tipbox}

\bigseparator

% --- 8.2.12 DPO ---
\subsubsection{DPO: Direct Preference Optimization}

\begin{defbox}[DPO]
\textbf{Problem with PPO:} Complex pipeline (generate $\to$ get reward $\to$ policy gradient $\to$ clip).

\textbf{DPO solution:} Directly optimise preferences without RL.

\textbf{Key idea:} We already know which responses humans prefer. Just make preferred responses more likely than dispreferred ones directly.

\textbf{Advantages over PPO:}
\begin{itemize}
  \item No reward model needed during training
  \item No reinforcement learning required
  \item Simpler and cheaper to implement
\end{itemize}
\end{defbox}

\bigseparator

% --- 8.2.13 GRPO ---
\subsubsection{GRPO: Group Relative Policy Optimization}

\begin{defbox}[GRPO]
\textbf{Problem with PPO:} Needs a critic/value function to estimate advantages.

\textbf{GRPO solution:} Compute advantages by comparing responses \textbf{within a group}.

\textbf{Process:}
\begin{enumerate}
  \item For a prompt, generate multiple candidate responses
  \item Compute statistics (mean, std) over the group
  \item Define advantage as how much better each response is relative to peers
\end{enumerate}

\textbf{Advantages:}
\begin{itemize}
  \item No learned value function needed (simpler)
  \item More stable training
  \item Reward function can be simple (e.g., response length)
\end{itemize}
\end{defbox}

\bigseparator

% --- Summary Table ---
\subsubsection{Summary: Training Methods Comparison}

\begin{center}
\renewcommand{\arraystretch}{1.3}
\begin{tabular}{@{} l p{4cm} p{4cm} @{}}
\toprule
\textbf{Method} & \textbf{Key Idea} & \textbf{Requirements} \\
\midrule
Pre-training & Next-token prediction & Plain text (no labels) \\
SFT & Train on instruction-response pairs & Labelled instruction data \\
RLHF + PPO & RL with clipped updates & Reward model + two LLMs in memory \\
LoRA & Low-rank weight decomposition & Reduced memory \\
DPO & Direct preference optimisation & Preference pairs (no RL) \\
GRPO & Group-relative advantages & Multiple responses per prompt \\
\bottomrule
\end{tabular}
\end{center}

\bigseparator

\subsubsection{Key Equations Summary}

\begin{thmbox}[Essential Formulas]
\textbf{Pre-training loss (cross-entropy):}
\[
\mathcal{L} = -\sum_{t} \log P(w_t \mid w_1, \ldots, w_{t-1})
\]

\textbf{PPO ratio and clipping:}
\[
r = \frac{P_{\text{current}}}{P_{\text{previous}}}, \quad L^{\text{CLIP}} = \min(r \cdot A, \text{clip}(r, 1-\varepsilon, 1+\varepsilon) \cdot A)
\]

\textbf{LoRA decomposition:}
\[
W' = W + L \times R, \quad \text{params} = m \times r + r \times n \ll m \times n
\]
\end{thmbox}

